% Latin-script Interslavic source component: Work 45 (Kapferer--Noether paper).
% Audited editorial provenance: ../00_SESSION_INDEPENDENT_STATE_v019.json and ../NORMALIZATION_DECISIONS_v019.jsonl.
\documentclass[11pt]{article}
\usepackage{fontspec}
\IfFontExistsTF{Times New Roman}{\setmainfont{Times New Roman}}{\setmainfont{TeX Gyre Termes}}
\IfFontExistsTF{Arial}{\setsansfont{Arial}}{\setsansfont{TeX Gyre Heros}}
\IfFontExistsTF{Consolas}{\setmonofont{Consolas}}{\setmonofont{Latin Modern Mono}}
\usepackage{amsmath,amssymb,mathtools}
\usepackage{amsthm}
\usepackage{geometry}
\usepackage{microtype}
\geometry{a4paper,margin=2.05cm}
\setlength{\parindent}{1.1em}
\setlength{\parskip}{0pt}
\newcommand{\foreign}[1]{#1}
\theoremstyle{plain}
\newtheorem{satz}{Teorem}
\newtheorem{zusatz}{Dodatok}
\emergencystretch=220em
\allowdisplaybreaks
\sloppy
\hbadness=10000
\vbadness=10000
\hfuzz=5pt
\vfuzz=12pt
\raggedbottom
\begin{document}

\newpage
\part*{Čest III\hspace{1em}Statja Kapferera --- Noether}
\addcontentsline{toc}{part}{Čest III: Statja Kapferera --- Noether}

\begin{center}
\Large\textbf{Potrěbne i dostatočne uslovja mnogokratnosti}\\[0.3em]
\Large\textbf{k Neterovomu fundamentalnomu teoremu}\\[0.5em]
\large\textbf{algebraičnyh funkcij}\\[1em]
\normalsize\textit{H. Kapferer}\\[0.5em]
\small\textit{(s dodatkom, společně s E. Noether)}\\[1em]
\normalsize \foreign{Math.\ Ann.}\ \textbf{97} (1927), S.~559--567
\end{center}

\vspace{1em}

Slědujuće prědstavja zaostrenje Neterovogo fundamentalnogo
teorema v formě potrěbnyh i dostatočnyh uslovij
mnogokratnosti\textsuperscript{1)}. Kompleks faktov Neterovogo
fundamentalnogo teorema ja, v oprěnju na idealno-teoretično
ponimanje teorema\textsuperscript{2)}, ale bez upotrěby
idealno-teoretičnyh pojmov i teoremov tu i v daljšem, svedu do
slědujućih dveh oddělnyh tvrđenij\textsuperscript{3)}:

\smallskip
\textbf{Teorem Ia.} Dla vsakoj obćej nulovej točky $x = a$,
$y = b$ dveh vzajemno prostyh polinomov $\varphi$, $\psi$ v $x$, $y$,
s libovoljnymi kompleksnymi koeficientami, egzistuje prinajmenje cělo
racionalno pozitivno čislo $\varrho$ tako, že modul
$(\varphi, \psi, \mathfrak{p}^{\varrho})$ označa točno tu samu
celost polinomov kako modul $(\varphi, \psi, \mathfrak{p}^{\sigma})$,
kde $\mathfrak{p} = (x-a, y-b)$ i $\sigma$ označa libovoljno cělo
racionalno čislo večše od $\varrho$\textsuperscript{4)}.

\smallskip
\textbf{Teorem Ib.} Ako se ta pojmovno jednoznačno oprěděljena
celost polinomov označi s $\mathfrak{q}$, oziroma s
$\mathfrak{q}_{\nu}$, kogda ide o $\nu$-toj iz $s$ različnyh nulevyh
toček, t.j. o $x = a_{\nu}$, $y = b_{\nu}$, togda imaje město:

Potrěbno i dostatočno za to, že dany polinom $K$ v $x$, $y$ imaje
svojstvo dělimości $K \equiv 0(\varphi, \psi)$, jest simultanno
izpolnjenje slědujućih $s$ kongruencij
\[
K \equiv 0(\mathfrak{q}_{1}), \quad K \equiv 0(\mathfrak{q}_{2}), \quad\ldots,\quad
K \equiv 0(\mathfrak{q}_{s}).
\]

Historično, sravni uvod osnovnoj raspravy Maksa Netera o toj
temě\textsuperscript{5)}, fundamentalny teorem vyrastl iz pokusov
rěšiti slědujuću fundamentalnu zadaču: pri vsakom danom polinomu
$K$ v $x$, $y$ trěba uměti ustanoviti, je-li on dělimy s
$(\varphi, \psi)$ ili ne, t.j. takože, egzistujeta-li dva druge
polinomy $A$ i $B$ v $x$, $y$ tako, že jednačenje
\[
K = A\varphi + B\psi
\]
stane identiteto v $x$, $y$.

Ta zadača ne rěšaje se Neterovym teoremom; poslědnji trěba bolje
razuměti samo kako, hoti fundamentalnu, redukciju\textsuperscript{6)}
postavjenoj zadači k prostějši, imenujuče k pytanju o potrěbnyh i
dostatočnyh uslovjah za obstojanje pojedinyh kongruencij
$K \equiv 0(\mathfrak{q})$. Kogda jest $K \equiv 0(\mathfrak{q})$?
Samo kogda to poslědnje pytanje jest odgovorjeno, jest takože
izpolnjena počętkova formulacija M.\ Netera. To jest cilj toj
raspravy, ktory takože faktično jest dosěgnuty, i to pokazovanjem
konečno mnogyh uslovij mnogokratnosti, praktično bez težkosti
kontrolovateljnyh, ktore za $K \equiv 0(\mathfrak{q})$ sut
potrěbne i dostatočne. Dostatočno uslovje za
$K \equiv 0(\mathfrak{q})$, ktore ale ne jest takože potrěbno
uslovje, jest poznano i često upotrěbjano i citovano v literaturi,
napriměr v formě:

\smallskip
\textit{Za obstojanje kongruencije $K \equiv 0(\varphi, \psi)$
dostatočno jest, že vsaka točka prěseka $\varphi = 0$, $\psi = 0$
v polinomu $K$ pojavja se s ``dostatočno'' vysokoju mnogokratnostju.}

\smallskip
Taj teorem sodrživaje se v teoremu I kako specialny slučaj; zato že
$\varrho$-kratna točka v $K$ jest ravnoznačna s
$K \equiv 0(\mathfrak{p}^{\varrho})$; togda a fortiori jest takože
$K \equiv 0(\varphi, \psi, \mathfrak{p}^{\varrho})$, t.j. takože
$K \equiv 0(\mathfrak{q})$. Ale obratno tvrđenje ne imaje město.

\vspace{0.5em}
\noindent\rule{\textwidth}{0.4pt}
\vspace{0.3em}

{\footnotesize
\textsuperscript{1)} Za podrobno prědstavjenje odkazuju na svoju rabotu s tym samym
naslovom v specialnom svazku sprav Heidelberžskoj akademije 1927
``\foreign{Beiträge zur Algebra}'', započetom \foreign{A.\ Loewy}, Freiburg i.\ Br., ktory
sodrživaje raboty raznyh avtorov.

\textsuperscript{2)} Sravni napr. \foreign{B.\ L.\ van der Waerden}, ``\foreign{Zur Nullstellentheorie der Polynomideale}''.
\foreign{Math.\ Annalen} 96 (1926), S.~203.

\textsuperscript{3)} V \S~1 citovanoj raspravy ja dokazal oba tvrđenja pojedino i direktno,
t.j. bez idealno-teoretičnyh pojmov.

\textsuperscript{4)} Bukvy $\mathfrak{p}$ i $\mathfrak{q}$ sut vybrane zato, aby napominati
o smyslovom soglasju označenogo s pojmami prosty ideal
$\mathfrak{p}$ i primarny ideal $\mathfrak{q}$, obyčnymi v novějšoj literaturi.

\textsuperscript{5)} \foreign{Max Noether}, ``\foreign{Über einen Satz aus der Theorie der algebraischen Funktionen}'',
\foreign{Math.\ Annalen} 6 (1872).

\textsuperscript{6)} Nazvana redukcija jest približno sravnima s toj, koju vsak zna iz
elementarnoj teorije čisel: svojstvo dělimości prirodnogo čisla
$A$ s takim samym $B$ jest ekvivalentno svojstvu dělimości $A$ s
$p_{1}^{e_{1}}$, s $p_{2}^{e_{2}}$, \ldots, s $p_{s}^{e_{s}}$, kde
$p_{1}, p_{2}, \ldots, p_{s}$ označajut različne proste dělitelji
$B = p_{1}^{e_{1}}p_{2}^{e_{2}} \cdots p_{s}^{e_{s}}$. Kako analog
stepena prostogo čisla trěba razuměti točno pojem ``primarny ideal'',
v tekstu označeny s $\mathfrak{q}$. [Sr. E.\ Noether,
``\foreign{Idealtheorie in Ringbereichen}'', \foreign{Math.\ Annalen} 83 (1921), uvod.]
}

\newpage
\section*{Glavny teorem}

Oběčano rěšenje problema možno kratko karakterizovati tako:
pojedinomu $\mathfrak{q}$ možno pridružiti sistem od $t$ pozitivnyh
cělyh racionalnyh čisel $(g_{1}), (g_{2}), \ldots, (g_{t})$, a
modulu $(K, \mathfrak{q})$ ravno toliko cělyh racionalnyh pozitivnyh
ili negativnyh čisel $(K_{1}), (K_{2}), \ldots, (K_{t})$, tako že
imaje město slědujući \textbf{glavny teorem}:

\begin{satz}[Glavny teorem]\label{satz:kap-II}
Potrěbno i dostatočno za $K \equiv 0(\mathfrak{q})$ jest
simultanno izpolnjenje $t$ uslovij
\[
(K_{1}) \geq (g_{1}), \quad (K_{2}) \geq (g_{2}), \quad\ldots,\quad (K_{t}) \geq (g_{t}).
\]
Pri tom
\[
(g_{1}) + (g_{2}) + \cdots + (g_{t}) = \mu,
\]
kde $\mu$ označa mnogokratnost prěseka, odgovarajuču
$\mathfrak{q}$ i oprěděljenu rezultantoju.
\end{satz}

V dokazu, bez ograničenja obće validnosti, prědpoložimo $x = 0$,
$y = 0$ kako nulevu točku para $(\varphi, \psi)$, odgovarajuču
$\mathfrak{q}$. Dalje, ako $A(x,y)$ jest kako-libo polinom\textsuperscript{7)}
v $x$, $y$, simbol $(A)$ bude označati čislo, ktore kaže, koliko
krat $y$ vystupa kako faktor v $A(0,y)$; ako ale $A(0,y)$
identično izčezne, togda $(A)$ stavjaje se ravno $\infty$. Čisla
$(g_{1}), (g_{2}), \ldots, (g_{t})$, pridružene k $\mathfrak{q}$,
oprěděljajut se slědujućim sistemom jednačenij, ktory tvori
metodičnu osnovu postupka:
\begin{equation}\label{eq:kap-1}
\begin{aligned}
f_{1}\cdot g_{1}(0,y):y^{(g_{1})} - g_{1}\cdot f_{1}(0,y):y^{(g_{1})} &= x\cdot h_{1},\\
f_{2}\cdot g_{2}(0,y):y^{(g_{2})} - g_{2}\cdot f_{2}(0,y):y^{(g_{2})} &= x\cdot h_{2},\\
&\vdots \qquad\qquad\vdots\qquad\qquad\vdots\\
f_{r}\cdot g_{r}(0,y):y^{(g_{r})} - g_{r}\cdot f_{r}(0,y):y^{(g_{r})} &= x\cdot h_{r}.
\end{aligned}
\end{equation}
On buduje se samo iz danogo para polinomov $\varphi$, $\psi$:
polinomy $f_{1}, g_{1}$ najprvo oprěděljajut se kako identične s
$\varphi, \psi$, eventualno po prěstavljanju poręda; porędok
dolžen byti tak, že $(f_{1}) \geq (g_{1})$. Pri toj fiksaciji, i
samo togda, $h_{1}$ bude cělo-racionalny.

Polinomy $f_{2}, g_{2}$ oprěděljajut se kako identične s
$h_{1}, g_{1}$ pri pobočnom uslovju $(f_{2}) \geq (g_{2})$. Obće par
polinomov $f_{r}, g_{r}$ oprědělja se kako identičny s
$h_{r-1}, g_{r-1}$ pri pobočnom uslovju $(f_{r}) \geq (g_{r})$.

Sistem jednačenij \eqref{eq:kap-1} možno istinno bezkončno
prodolžati v objasnjenom sposobu; glavny teorem ale potrěbuje
stvoriti samo prve $t$ jednačenja sistema, kde $t$ jest oprěděljeno
slědujućim faktom:

\smallskip
\textbf{Pomočny teorem I.} Egzistuje prirodno čislo $t$ [i to
$t \leq \mu$, pod $\mu$ razuměvaje se mnogokratnost prěseka, v kojoj
točka $x = 0$, $y = 0$ vystupa v $\varphi = 0$, $\psi = 0$] tako, že
prve $t$ čisla, oprěděljene s \eqref{eq:kap-1}, imenujuče
$(g_{1}), (g_{2}), \ldots, (g_{t})$, sut vsě večše od nule, pokym
jednočasno $(g_{t+1})$ i vsě slědujuće sut ravne nuli.

S urědnjenjem produktnogo teorema rezultantov, mnogokratnost
prěseka v točce $x = 0$, $y = 0$ za $f_{r} = 0$, $g_{r} = 0$ bude
ravna
\[
\mu - (g_{1}) - (g_{2}) - \cdots - (g_{r}),
\]
iz čego slěduje prěryv po najvyše $\mu$ krokah. Iz toga jednočasno
slěduje fakt
\begin{equation}\label{eq:kap-2}
\mu = (g_{1}) + (g_{2}) + \cdots + (g_{t}),
\end{equation}
formula značitelna zato, že ona obće daje racionalnu i praktično
legko izvršimu metodu oprěděljenja mnogokratnosti prěseka, kogda
odgovarajuča nuleva točka jest poznana.

Kogda čisla $(g_{i})$ sut tako poznane, čisla $(K_{i})$ oprěděljajut
se slědujućim sistemom jednačenij, odgovarajučim k \eqref{eq:kap-1}:
\begin{equation}\label{eq:kap-3}
\begin{aligned}
K_{1}\cdot g_{1}(0,y):y^{(g_{1})} - g_{1}\cdot K_{1}(0,y):y^{(g_{1})} &= x\cdot K_{2},\\
K_{2}\cdot g_{2}(0,y):y^{(g_{2})} - g_{2}\cdot K_{2}(0,y):y^{(g_{2})} &= x\cdot K_{3},\\
&\vdots \qquad\qquad\vdots\qquad\qquad\vdots\\
K_{r}\cdot g_{r}(0,y):y^{(g_{r})} - g_{r}\cdot K_{r}(0,y):y^{(g_{r})} &= x\cdot K_{r+1}.
\end{aligned}
\end{equation}
$K_{1}$ oprědělja se kako identičny s $K$. Jasno jest, že $K_{2}$
bude cělo-racionalny togda i samo togda, kogda
$(K_{1}) \geq (g_{1})$. Ako to poslědnje jest izpolnjeno, dalje
slěduje: $K_{3}$ bude cělo-racionalny togda i samo togda, kogda
$(K_{2}) \geq (g_{2})$. Obće, ako $K_{r}$ už jest cělo-racionalny,
$K_{r+1}$ bude cělo-racionalny togda i samo togda, kogda
$(K_{r}) \geq (g_{r})$.

Dokaz glavnogo teorema (teorema~\ref{satz:kap-II}) teraz oprě se na
daljši

\smallskip
\textbf{Pomočny teorem II.} Ako se
$(f_{i}, g_{i}, \mathfrak{p}^{e})$ postavi ravno
$\mathfrak{q}^{(i)}$, togda za vsako $i$ iz cěločiselnoj serije
$1, 2, \ldots$ \textit{in inf.} imaje město teorem:

Za kongruenciju
\[
K_{i} \equiv 0(\mathfrak{q}^{(i)})
\]
potrěbno i dostatočno jest simultanno izpolnjenje dveh uslovij:
\[
(K_{i}) \geq (g_{i}) \qquad \text{i} \qquad K_{i+1} \equiv 0(\mathfrak{q}^{(i+1)}).
\]

Dokaz\textsuperscript{8)} pomočnoho teorema dobyva se kombinovanjem
sistemov jednačenij \eqref{eq:kap-1} i \eqref{eq:kap-3}. Ako se
urědni, že $g_{t+1}$ po pomočnom teoremu I v nulevoj točce veče ne
izčezne, t.j. že modul $\mathfrak{q}^{(t+1)}$ sostoji se iz vsih
polinomov (jest jediničny ideal), i zato
$K_{t+1} \equiv 0(\mathfrak{q}^{(t+1)})$ stane trivialno, togda
$t$-kratno upotrěbjenje pomočnoho teorema II daje dokaz glavnogo
teorema.

\smallskip
Bez dokazov, ktore nahodet se v mojej gore citovanoj raspravi,
dodaju ješče dva dodatka k glavnomu teoremu:

\begin{zusatz}[I]
Čislo $t$ glavnogo teorema jest identično s prinajmenje\textsuperscript{9)}
pozitivnym čislom $\nu$ takoho roda, že
$x^{\nu} \equiv 0(\mathfrak{q})$.
\end{zusatz}

\begin{zusatz}[II]
Sistemy jednačenij \eqref{eq:kap-1} i \eqref{eq:kap-3}, i te, ktore
iz njih vznikajut prěstavljanjem $x$ s $y$, pri vhodnoj kombinaciji
davajut racionalnu metodu točnogo čislovnogo oprěděljenja
``naležečego'' k $\mathfrak{q}$ eksponenta
$\varrho$, t.j. togo minimalnogo eksponenta $\varrho$ iz teorema Ia.
\end{zusatz}

E.\ Bertini\textsuperscript{10)} dal obću formulu za $\varrho$.
Že ta formula ale ne vsegda daje minimalnu vrědnost $\varrho$,
pokazano jest v citovanoj raspravi na prikladu
$\varphi = x^{3} + y^{4}$, $\psi = x^{2} + y^{3}$. Nakonec trěba
ješče zamětiti, že i sama formula Bertinija može byti po novomu
osnovana našim glavnym teoremom.

\vspace{0.5em}
\noindent\rule{\textwidth}{0.4pt}
\vspace{0.3em}

{\footnotesize
\textsuperscript{7)} Ako za $A(x,y)$ dopustimo takože drobno-racionalne
funkcije $C(x,y):D(x,y)$, togda $(A) = (C) - (D)$; v poslědnjem
slučaju može $(A)$ iměti i negativne vrědnosti.

\textsuperscript{8)} Sr. takože dokaz v zamětce I ``Dodatka''.

\textsuperscript{9)} Tu se zato dobyvajut točne eksponenty v protivnosti k
ocěnkam gospođice Grete Hermann [``\foreign{Die Frage der endlich vielen Schritte
in der Theorie der Polynomideale}'' (s upotrěboju posmrtno ostavjenyh
teoremov K.\ Hentzelt), \foreign{Math.\ Annalen} 95 (1926)], kde pri
mnogo obćejšoj postavjenosti problema davajut se samo gorne granice.

\textsuperscript{10)} E.\ Bertini, \foreign{Math.\ Annalen} 34 (1889).
}

\newpage
\section*{Teorem III i dodatok}

Interesna specializacija glavnogo teorema slěduje iz ograničajučego
prědpoloženja, že nuleva točka odgovarajučego
$\mathfrak{q} = (\varphi, \psi, \mathfrak{p}^{e})$ jest prosta
točka prinajmenje v jednoj iz dveh krivulj $\varphi = 0$, $\psi = 0$.
Ako ona jest napr. prosta točka v $\psi$, togda možno pokazati, že
čislo $t$ glavnogo teorema jest točno ravno $\mu$ i že $t$
neravnosti glavnogo teorema možno zaměniti jedino uslovjem
mnogokratnosti. To vodi k slědujućemu obćemu teoremu:

\begin{satz}[Teorem III]\label{satz:kap-III}
Ako vsě točky prěseka $\varphi = 0$, $\psi = 0$ sut proste točky v
$\psi$, togda, prědpoloživši že $K$ jest vzajemno prosty s $\psi$,
za obstojanje kongruencije $K \equiv 0(\varphi, \psi)$ potrěbno i
dostatočno jest, že $K$ izčezaje v vsih tyh točkah, ale tako, že
mnogokratnost prěseka v paru krivulj $K = 0$, $\psi = 0$ v vsakoj
iz njih jest prinajmenje taka velika kako v paru krivulj
$\varphi = 0$, $\psi = 0$.
\end{satz}

Taj teorem III jest ješče iz osobnoj pričiny značitelny: uslovje
mnogokratnosti, ktore on potrěbuje, jest tako, že jego izpolnjenje
ili neizpolnjenje može byti ustanovjeno konečno-kratnym
diferencovanjem, nezavisno od formuly \eqref{eq:kap-2}, oprětoj na
\eqref{eq:kap-1}. To jest neposrědnje jasno, ako cituju slědujući
``teorem'', ktory ja postavil v 1923 roku\textsuperscript{11)}:

\smallskip
\textit{Ako $A$ i $B$ sut kako-libo dva vzajemno proste polinomy v
$x$, $y$, a $P$ jest obća točka, ktora jest prosta točka od $B$,
togda mnogokratnost toj točky kako točky prěseka jest togda i samo
togda ravna čislu $\mu$, kogda mnogokratnost toj samej točky v paru
krivulj $A' = 0$, $B = 0$ jest točno $\mu - 1$; pri tom označa}
\[
A' = \frac{\partial A}{\partial x} \cdot \frac{\partial B}{\partial y} -
\frac{\partial A}{\partial y} \cdot \frac{\partial B}{\partial x}.
\]

Nakonec trěba ješče spomenuti, že glavny teorem
(teorem~\ref{satz:kap-II}) s maloju modifikacijoju može byti
prěneseny takože na homogene polinomy v 3 variablah $x$, $y$, $z$.
To prěnesenje udaje se osobito elegantno v gore navednom specialnom
slučaju. Imenno možno pokazati:

\begin{zusatz}[k teoremu III]
Teorem~\ref{satz:kap-III} ostavaje cělkom nezměnjeny, i po slovah,
ako pod $\varphi$, $\psi$, $K$ razuměvajut se homogene polinomy v
treh variablah\textsuperscript{12)}.
\end{zusatz}

\vspace{0.5em}
\noindent\rule{\textwidth}{0.4pt}
\vspace{0.3em}

{\footnotesize
\textsuperscript{11)} H.\ Kapferer, ``\foreign{Über die Multiplizität der Schnittpunkte von zwei algebraischen
Kurven}''. \foreign{Jahresber.\ der D.\ M.-V.} 32 (1923).

\textsuperscript{12)} Pytanje o uslovjah za obstojanje homogene kongruencije v
treh variablah $K(x,y,z) \equiv 0(\varphi(x,y,z), \psi(x,y,z))$ jest,
hoti pod silno specializovanymi prědpoloženjami, takože jadro
izslědovanja \foreign{A.\ Ostrowski}, ``\foreign{Über die Maxwellsche Erzeugung der
Kugelfunktionen}'', \foreign{Jahresb.\ der D.M.-V.} 33 (1925). V toj
notě prědpoloženo jest: $\varphi$ razpadaje se na same linearne
faktory; $\varphi$ i $K$ sut jednakoho poręda; $\psi = x^{2} + y^{2} + z^{2}$.
Zato že tu $\psi$ jest cělkom bez singularnostij, glavny teorem
Ostrovskijevoj noty možno razuměti kako specialny slučaj našego
teorema III.
}

\newpage
\section*{Dodatok, společně s E. Noether\textsuperscript{13)}}

Sistem jednačenij \eqref{eq:kap-1} daje jednočasno racionalnu
metodu za oprěděljenje polnogo sistema ostatkov po $\mathfrak{q}$,
pod čim trěba razuměti polny sistem prědstaviteljev konečno mnogyh,
v odnosu k oblasti koeficientov $P$, modulo $\mathfrak{q}$ linearno
nezavisnyh ostatkovyh klasov. S tym pojedine kroky v dokazu
glavnogo teorema stavajut mnogo prozrěčne; jednočasno kako
pobočny rezultat vyhodi fakt, že rezultanta-mnogokratnost $\mu$
soglasuje se s dolžinoju ideala $\mathfrak{q}$;

\vspace{0.5em}
\noindent\rule{\textwidth}{0.4pt}
\vspace{0.3em}

{\footnotesize
\textsuperscript{13)} Idealno-teoretično tolkovanje pohodit od E.\ Noether, a osnovny
dokaz (5) i (6) od H.\ Kapferer.

\textsuperscript{14)} Polje $P$ trěba, bez ograničenja obće validnosti, prědpoložiti
algebraično zatvorenym; v tom obsegu faktično imajut město vsě poprědne
teoremy, s izključenjem diferenciacijskogo uslovja. Dolžina
$\mathfrak{q}$ jest takože oprěděljena kako dolžina kompozicijnoj
serije idealov, iduče od $\mathfrak{p}$ do $\mathfrak{q}$. Sr.
E.\ Noether, \foreign{Jahresber.\ d.\ d.\ Math.-Ver.} 34 (1925), S.~101
(kosa paginacija); podrobněje u H.\ Grell, \foreign{Math.\ Annalen} 97
(1927), S.~529.

Dokazy v literaturi za soglasje obojih čisel mnogokratnosti (pri
$n$ variablah) sut zelo složne; najprostějši jest ješče nepublikovany
posmrtno ostavjeny dokaz Hentzelta.
}

\vspace{1em}
t.j. s rangom prstena ostatkovyh klasov ili s čislom v odnosu k $P$
linearno nezavisnyh ostatkovyh klasov\textsuperscript{14)}.

Polny sistem ostatkov jest, s označenjami iz \eqref{eq:kap-1} i
pomočnoho teorema II, karakterizovany slědujućim

\begin{satz}[Teorem IV]\label{satz:kap-IV}
Polny sistem ostatkov po $\mathfrak{q}$ dany jest polinomami
\[
h_{t},\ h_{t}y,\ldots,h_{t}y^{(g_{t})-1};\ \ldots;\
h_{i},\ h_{i}y,\ldots,h_{i}y^{(g_{i})-1};\ \ldots;\
h_{1},\ h_{1}y,\ldots,h_{1}y^{(g_{1})-1}.
\]
Pri tom vsaky raz
$h_{i}, h_{i}y, \ldots, h_{i}y^{(g_{i})-1}$ daje polny
sistem ostatkov od $\mathfrak{q}^{(i+1)}$ do
$\mathfrak{q}^{(i)}$.
\end{satz}

Iz definicije slěduje neposrědnje:
$\mathfrak{q}^{(i+1)} = (\mathfrak{q}^{(i)}, h_{i})$, zato
$\mathfrak{q}^{(i)} \equiv 0(\mathfrak{q}^{(i+1)})$.
Teorem~\ref{satz:kap-IV} bude dokazany, ako pokazano jest:
\begin{equation}\label{eq:kap-4}
h_{i}x \equiv 0(\mathfrak{q}^{(i)}); \qquad
h_{i}y^{(g_{i})} \equiv 0(\mathfrak{q}^{(i)})
\end{equation}
i
\begin{equation}\label{eq:kap-5}
\text{iz } h_{i}F(y) \equiv 0(\mathfrak{q}^{(i)})
\quad \text{slěduje} \quad F(y) \equiv 0(y^{(g_{i})}).
\end{equation}

Relacije \eqref{eq:kap-4} imenno govoret, že linearne kombinacije
$h_{i}, h_{i}y, \ldots, h_{i}y^{(g_{i})-1}$ izčerpavajut
ostatkove klasy od $\mathfrak{q}^{(i+1)}$ do
$\mathfrak{q}^{(i)}$, pokym \eqref{eq:kap-5} vyražaje linearnu
nezavisnost. Zato že $\mathfrak{q}^{(t+1)}$ stavaje jediničnym
idealom; to teraz slěduje iz konečnoj dolžiny $\mathfrak{q}$,
iz toga dobyva se dokaz teorema~\ref{satz:kap-IV} i jednočasno
fakt, že dolžina $\mathfrak{q}$ stavaje ravna
$(g_{1}) + \cdots + (g_{t})$, zato po \eqref{eq:kap-2} soglasuje se
s rezultanta-mnogokratnostju.

Prva iz relacij \eqref{eq:kap-4} čita se iz \eqref{eq:kap-1};
druga dobyva se slědujućim sposobom:

Imamo $h_{i}g_{i} \equiv 0(\mathfrak{q}^{(i)})$; zato zaradi prve
relacije \eqref{eq:kap-4} takože
$h_{i}g_{i}(0,y) \equiv 0(\mathfrak{q}^{(i)})$, ili
$h_{i}y^{(g_{i})}(g_{i}(0,y):y^{(g_{i})}) \equiv 0(\mathfrak{q}^{(i)})$,
iz čego $h_{i}y^{(g_{i})} \equiv 0(\mathfrak{q}^{(i)})$
slěduje zaradi $g_{i}(0,y):y^{(g_{i})} \not\equiv 0(y)$.

Pred dokazom \eqref{eq:kap-5} dajmo opazku neposrědnje čitanu iz
\eqref{eq:kap-1}: $f_{i}$ i $g_{i}$ mogut iměti kako obći dělitelj
samo polinom $d_{i}(y)$ v $y$, ne dělimy s $y$:
$f_{i} = d_{i}(y)\bar{f}_{i}$, $g_{i} = d_{i}(y)\bar{g}_{i}$, i
$\bar{f}_{i}$ i $\bar{g}_{i}$ sut vzajemno proste. Zaradi
$d_{i}(y) \not\equiv 0(\mathfrak{p})$ bude $\mathfrak{q}^{(i)}$
takože ravno $(\bar{f}_{i}, \bar{g}_{i}, \mathfrak{p}^{e})$;
pri tom to bude ravno
$(\bar{f}_{i}, \bar{g}_{i}, \mathfrak{p}^{\sigma})$ za vsako
$\sigma \geq \varrho$, kako dělitelj $\mathfrak{q}$ za vsako
tako $\sigma$; zato $\mathfrak{q}^{(i)}$ jest primarna komponenta
$(\bar{f}_{i}, \bar{g}_{i})$, naležeča k nulevoj točce.

Nehaj $\gamma_{1}, \delta_{1}; \ldots; \gamma_{r}, \delta_{r}$ sut
nuleve točky $(\bar{f}_{i}, \bar{g}_{i})$, različne od
$x = 0$, $y = 0$, i nehaj $\tau_{1}, \ldots, \tau_{r}$ sut eksponenty
odgovarajučih primarnyh komponentov $(\bar{f}_{i}, \bar{g}_{i})$;
dalje v produktu
\[
P = ((x-\gamma_{1}) + u(y-\delta_{1}))^{\tau_{1}} \cdots
((x-\gamma_{r}) + u(y-\delta_{r}))^{\tau_{r}}
\]
neizvěstnu $u$ vyberemo kako element iz $P$ tako, že ostane
$P \not\equiv 0(\mathfrak{p})$. Ako se $G(x,y)$ postavi ravno
produktu specializovanogo $P$ s $d_{i}(y)$, togda
$G(x,y) \not\equiv 0(\mathfrak{p})$ i
$\mathfrak{q}^{(i)} \cdot G(x,y) \equiv 0(f_{i}, g_{i})$.

Iz prědpoloženja $h_{i}F(y) \equiv 0(\mathfrak{q}^{(i)})$
slěduje
$h_{i}F(y)G(x,y) \equiv 0(f_{i}, g_{i})$, ili, spoločno
s \eqref{eq:kap-1}, v formě identitetov:
\[
h_{i}F(y)G = Af_{i} + Bg_{i}; \qquad
h_{i}x = (g_{i}(0,y):y^{(g_{i})})\cdot f_{i}
- (f_{i}(0,y):y^{(g_{i})})g_{i}.
\]
To daje ale
\[
\bar{f}_{i}(xA - F(y)H) \equiv 0(\bar{g}_{i}); \qquad
H = G(x,y) \cdot g_{i}(0,y) : y^{(g_{i})} \not\equiv 0(\mathfrak{p}).
\]
Iz vzajemnoj prostosti $\bar{f}_{i}$ i $\bar{g}_{i}$ slěduje dalje
\[
xA - F(y)H \equiv 0(\bar{g}_{i});
\]
tako $F(y)H \equiv 0(\bar{g}_{i}, x)
\equiv 0(\bar{g}_{i}, x, \mathfrak{p}^{\sigma})
= (y^{(g_{i})}, x)$, i s tym $F(y) \equiv 0(y^{(g_{i})})$.
Tym jest \eqref{eq:kap-5}, i zato teorem~\ref{satz:kap-IV}, dokazany.

\smallskip
\textbf{Zamětka I.} Relacije \eqref{eq:kap-4} i \eqref{eq:kap-5}
spoločno s pomočnym teoremom II možno zapisati kako vzajemno
recipročne relacije med kvocientami idealov\textsuperscript{15)};
pri tom dokaz drugoj relacije daje jednočasno dokaz pomočnoho
teorema II:
\begin{equation}\label{eq:kap-6}
\mathfrak{q}^{(i)} : \mathfrak{q}^{(i+1)}
= (\mathfrak{q}^{(i)}, x) = (y^{(g_{i})}, x); \qquad
\mathfrak{q}^{(i)} : (\mathfrak{q}^{(i)}, x)
= \mathfrak{q}^{(i+1)}.\qquad\text{\textsuperscript{16)}}
\end{equation}
Prva relacija prědstavja sjedinenje \eqref{eq:kap-4} i
\eqref{eq:kap-5}; k drugoj trěba zamětiti: v \eqref{eq:kap-4}
bylo dokazano
$x \cdot \mathfrak{q}^{(i+1)} \equiv 0(\mathfrak{q}^{(i)})$;
obratno dobyva se analogično kako \eqref{eq:kap-5}. Iz
$xM \equiv 0(\mathfrak{q}^{(i)})$ slěduje kako vyše:
\[
xMG \equiv 0(f_{i}, g_{i}); \quad \text{tako }\quad
xMG \cdot g_{i}(0,y) : y^{(g_{i})}
\equiv 0(xh_{i}, g_{i}).
\]
Zato že $g_{i} \not\equiv 0(x)$ po definiciji, imenujuče
$(f_{i}) \geq (g_{i})$, iz toga vyhodi
\[
M \cdot (G \cdot g_{i}(0,y) : y^{(g_{i})})
\equiv 0(h_{i}, g_{i}) \equiv 0(\mathfrak{q}^{(i+1)}),
\quad\text{i zato}\quad M \equiv 0(\mathfrak{q}^{(i+1)}).
\]

Že uslovje pomočnoho teorema II jest potrěbno, slěduje teraz tako:
$K_{i} \equiv 0(\mathfrak{q}^{(i)})$ daje
$K_{i}(0,y) \equiv 0(y^{(g_{i})}, x)$; zato
$(K_{i}) \geq (g_{i})$ i s tym, po \eqref{eq:kap-3},
\[
x \cdot K_{i+1} \equiv 0(\mathfrak{q}^{(i)});
\]
zato po drugoj relaciji \eqref{eq:kap-6} bude
$K_{i+1} \equiv 0(\mathfrak{q}^{(i+1)})$.

Že uslovje jest dostatočno, čita se neposrědnje; konečno-kratno
upotrěbjenje pomočnoho teorema II daje ale glavny teorem. Takože
iteracija drugoj relacije \eqref{eq:kap-6} daje dodatok I, imenujuče
$\mathfrak{q} : x^{t} = \mathfrak{q}^{(t+1)}$, zato jediničny
ideal, i $t$ jest najniži tak stepen, zaradi
$\mathfrak{q} : x^{t-1} = \mathfrak{q}^{(t)}$.

\textbf{Zamětka II.} Za polinom $K$, ne dělimy s $\mathfrak{q}$,
možno ješče tvrditi: ako
$K \equiv 0(\mathfrak{q}^{(i+1)})$,
$\not\equiv 0(\mathfrak{q}^{(i)})$, togda
$(K, \mathfrak{q}^{(i)}) = (h_{i}y^{(K)}, \mathfrak{q}^{(i)})$.
Zato že po teoremu~\ref{satz:kap-IV}
$K \equiv h_{i}y^{\lambda}c(y)$ modulo
$(\mathfrak{q}^{(i)})$ s $c(y) \not\equiv 0(y)$; tako vyhodi
\[
(K, \mathfrak{q}^{(i)}) = (h_{i}y^{\lambda}, \mathfrak{q}^{(i)}).
\]
Pri tom $y^{(g_{i})-\lambda} \cdot K \equiv 0(\mathfrak{q}^{(i)})$,
i to po \eqref{eq:kap-5} jest najniži tako stepen. Po pomočnom
teoremu II jest ale $(g_{i}) - (K)$ toj najniži stepen; zato
$\lambda$ stavaje ravno $(K)$.

\vfill
\begin{flushright}
(Prišlo 24.\ 10.\ 1926.)
\end{flushright}

\vspace{0.5em}
\noindent\rule{\textwidth}{0.4pt}
\vspace{0.3em}

{\footnotesize
\textsuperscript{15)} Pod kvocientom $\mathfrak{a}:\mathfrak{b}$, kako poznano,
razuměje se celost polinomov $c(x,y)$ takyh, že
$\mathfrak{b}c \equiv 0(\mathfrak{a})$.

\textsuperscript{16)} Iz tyh dveh vzajemno recipročnyh relacij jedna uslovjuje
drugu po obćej teoriji ireducibilnyh idealov; sr. Macaulay,
\textit{\foreign{Modular Systems}}, Nr.\,72 i 73. \foreign{Cambridge Tracts} 19 (1916).
}

\end{document}
